\section*{2.6 基和秩}
\addcontentsline{toc}{section}{2.6 基和秩}
\label{sec:2.6}

在向量空间 $ V $ 中，我们对这样一类向量集合 $ A $ 尤其感兴趣：它们具备使得 $ V $ 中的任意向量 $ \boldsymbol{v} $ 均可通过 $ A $ 中向量的线性组合来获得的性质。这些向量属于特殊向量，在下文中我们将对其性质进行刻画。

\section*{2.6.1 生成集与基}
\addcontentsline{toc}{subsection}{2.6.1 生成集与基}

\textbf{定义 2.13}（生成集与张成空间）。考虑一个向量空间 $ V = (V, +, \cdot) $ 以及向量集合 $ A = \{\boldsymbol{x}_1, \ldots, \boldsymbol{x}_k\} \subseteq V $。如果 $ V $ 中的每一个向量 $ \boldsymbol{v} $ 都可以表示为 $ \boldsymbol{x}_1, \ldots, \boldsymbol{x}_k $ 的线性组合，则称 $ A $ 为 $ V $ 的 \emph{生成集（generating set）} 。由 $ A $ 中向量的所有线性组合构成的集合称为 $ A $ 的 \emph{张成空间（span）} 。如果 $ A $ 张成了向量空间 $ V $，我们记作 $ V = \operatorname{span}[A] $ 或 $ V = \operatorname{span}[\boldsymbol{x}_1, \ldots, \boldsymbol{x}_k] $。

生成集是张成整个向量（子）空间的向量集合，即该空间中的每个向量都可以表示为生成集中向量的线性组合。现在，我们将更加具体地刻画张成一个向量（子）空间的“最小”生成集。

\textbf{定义 2.14}（基）。考虑一个向量空间 $ V = (V, +, \cdot) $ 及其子集 $ A \subseteq V $。如果不存在能够张成 $ V $ 且比 $ A $ 更小的集合 $ \tilde{A} \subsetneq A \subseteq V $，则称 $ V $ 的生成集 $ A $ 是 \emph{极小的（minimal）} 。$ V $ 的每一个线性无关的生成集都是极小的，并称为 $ V $ 的一组 \emph{基（basis）} 。

设 $ V = (V, +, \cdot) $ 为一个向量空间，且 $ B \subseteq V, B \neq \emptyset $。那么，以下陈述是等价的：

\begin{itemize}
\item $ B $ 是 $ V $ 的一组基。
\item $ B $ 是一个极小生成集。
\item $ B $ 是 $ V $ 中的极大线性无关向量集，即向该集合中添加任何其他向量都会使其变为线性相关的。
\item $ V $ 中的每一个向量 $ \boldsymbol{x} $ 都是来自 $ B $ 中向量的线性组合，且每个线性组合都是唯一的，即由： 且 $ \lambda_i, \psi_i \in \mathbb{R}, \boldsymbol{b}_i \in B $，可推导出对所有 $ i = 1, \ldots, k $ 均有 $ \lambda_i = \psi_i $。
\end{itemize}
\begin{equation}
\boldsymbol{x} = \sum_{i=1}^k \lambda_i \boldsymbol{b}_i = \sum_{i=1}^k \psi_i \boldsymbol{b}_i
\tag{2.77}\label{eq:2-77}
\end{equation}
\begin{itemize}
\end{itemize}

\begin{notebox}
一组基既是一个极小生成集，也是一个极大线性无关向量集合。
\end{notebox}

\begin{notebox}
\textbf{例 2.16}
\end{notebox}

\begin{notebox}
在 $ \mathbb{R}^3 $ 中， \emph{标准基（canonical/standard basis）} 为：

\begin{equation}
B = \left\{ \begin{bmatrix} 1 \\ 0 \\ 0 \end{bmatrix}, \begin{bmatrix} 0 \\ 1 \\ 0 \end{bmatrix}, \begin{bmatrix} 0 \\ 0 \\ 1 \end{bmatrix} \right\}
\tag{2.78}\label{eq:2-78}
\end{equation}

$ \mathbb{R}^3 $ 中的其他不同基例如：

\begin{equation}
\begin{adjustbox}{max width=\linewidth}
B_1 = \left\{ \begin{bmatrix} 1 \\ 0 \\ 0 \end{bmatrix}, \begin{bmatrix} 1 \\ 1 \\ 0 \end{bmatrix}, \begin{bmatrix} 1 \\ 1 \\ 1 \end{bmatrix} \right\}, \quad B_2 = \left\{ \begin{bmatrix} 0.5 \\ 0.8 \\ 0.4 \end{bmatrix}, \begin{bmatrix} 1.8 \\ 0.3 \\ 0.3 \end{bmatrix}, \begin{bmatrix} -2.2 \\ -1.3 \\ 3.5 \end{bmatrix} \right\}
\end{adjustbox}
\tag{2.79}\label{eq:2-79}
\end{equation}

而集合：

\begin{equation}
A = \left\{ \begin{bmatrix} 1 \\ 2 \\ 3 \\ 4 \end{bmatrix}, \begin{bmatrix} 2 \\ -1 \\ 0 \\ 2 \end{bmatrix}, \begin{bmatrix} 1 \\ 1 \\ 0 \\ -4 \end{bmatrix} \right\}
\tag{2.80}\label{eq:2-80}
\end{equation}

是线性无关的，但它不是 $ \mathbb{R}^4 $ 的生成集（也不是基）：例如向量 $ [1, 0, 0, 0]^\top $ 无法由 $ A $ 中元素的线性组合获得。
\end{notebox}

\emph{备注} 。每个向量空间 $ V $ 都拥有一组基 $ B $。上述例子表明一个向量空间 $ V $ 可以存在很多组基，即基并不是唯一的。然而，所有的基都拥有相同数量的元素，即 \emph{基向量（basis vectors）} 。

我们只考虑有限维向量空间 $ V $。在这种情况下，$ V $ 的 \emph{维度（dimension）} 等于 $ V $ 的基向量数量，记为 $ \operatorname{dim}(V) $。如果 $ U \subseteq V $ 是 $ V $ 的一个子空间，则 $ \operatorname{dim}(U) \leqslant \operatorname{dim}(V) $，且当且仅当 $ U = V $ 时 $ \operatorname{dim}(U) = \operatorname{dim}(V) $。直观地，向量空间的维度可以看作该向量空间中独立方向的数量。

\begin{notebox}
向量空间的维度对应于其基向量的数量。
\end{notebox}

\emph{备注} 。向量空间的维度并不必然等于向量中元素的个数。例如，向量空间

\begin{equation*}
V = \operatorname{span}\left[\begin{bmatrix} 0 \\ 1 \end{bmatrix}\right]
\end{equation*}

是一维的，尽管其基向量包含两个元素。

\emph{备注} 。求由向量张成的子空间 $ U = \operatorname{span}[\boldsymbol{x}_1, \ldots, \boldsymbol{x}_m] \subseteq \mathbb{R}^n $ 的一组基，可以通过执行以下步骤完成：

\begin{enumerate}
\item 将张成向量作为矩阵 $ \boldsymbol{A} $ 的列写入；
\item 确定 $ \boldsymbol{A} $ 的行阶梯形矩阵；
\item 与主元列相关联的原张成向量即构成 $ U $ 的一组基。
\end{enumerate}

\begin{notebox}
\textbf{例 2.17（确定一组基）}
\end{notebox}

\begin{notebox}
对于由以下向量张成的向量子空间 $ U \subseteq \mathbb{R}^5 $：

\begin{equation}
\begin{adjustbox}{max width=\linewidth}
\boldsymbol{x}_1 = \begin{bmatrix} 1 \\ 2 \\ -1 \\ -1 \\ -1 \end{bmatrix}, \quad \boldsymbol{x}_2 = \begin{bmatrix} 2 \\ -1 \\ 1 \\ 2 \\ -2 \end{bmatrix}, \quad \boldsymbol{x}_3 = \begin{bmatrix} 3 \\ -4 \\ 3 \\ 5 \\ -3 \end{bmatrix}, \quad \boldsymbol{x}_4 = \begin{bmatrix} -1 \\ 8 \\ -5 \\ -6 \\ 1 \end{bmatrix} \in \mathbb{R}^5
\end{adjustbox}
\tag{2.81}\label{eq:2-81}
\end{equation}

我们想要找出向量 $ \boldsymbol{x}_1, \ldots, \boldsymbol{x}_4 $ 中哪些构成了 $ U $ 的一组基。为此，我们需要检验 $ \boldsymbol{x}_1, \ldots, \boldsymbol{x}_4 $ 是否线性无关。因此，我们需要求解：

\begin{equation}
\sum_{i=1}^4 \lambda_i \boldsymbol{x}_i = \boldsymbol{0}
\tag{2.82}\label{eq:2-82}
\end{equation}

这导出了以如下矩阵为系数矩阵的齐次线性方程组：

\begin{equation}
\begin{adjustbox}{max width=\linewidth}
[\boldsymbol{x}_1, \boldsymbol{x}_2, \boldsymbol{x}_3, \boldsymbol{x}_4] = \begin{bmatrix} 1 & 2 & 3 & -1 \\ 2 & -1 & -4 & 8 \\ -1 & 1 & 3 & -5 \\ -1 & 2 & 5 & -6 \\ -1 & -2 & -3 & 1 \end{bmatrix}
\end{adjustbox}
\tag{2.83}\label{eq:2-83}
\end{equation}

利用线性方程组的初等变换规则，我们得到行阶梯形矩阵：

\begin{equation*}
\begin{adjustbox}{max width=\linewidth}
\begin{bmatrix} 1 & 2 & 3 & -1 \\ 2 & -1 & -4 & 8 \\ -1 & 1 & 3 & -5 \\ -1 & 2 & 5 & -6 \\ -1 & -2 & -3 & 1 \end{bmatrix} \leadsto \cdots \leadsto \begin{bmatrix} 1 & 2 & 3 & -1 \\ 0 & 1 & 2 & -2 \\ 0 & 0 & 0 & 1 \\ 0 & 0 & 0 & 0 \\ 0 & 0 & 0 & 0 \end{bmatrix}
\end{adjustbox}
\end{equation*}

由于主元列指示了哪些向量构成线性无关集，从行阶梯形中我们可以看出第 1、2、4 列为主元列，因此 $ \boldsymbol{x}_1, \boldsymbol{x}_2, \boldsymbol{x}_4 $ 是线性无关的（因为线性方程组 $ \lambda_1 \boldsymbol{x}_1 + \lambda_2 \boldsymbol{x}_2 + \lambda_4 \boldsymbol{x}_4 = \boldsymbol{0} $ 只有平凡解 $ \lambda_1 = \lambda_2 = \lambda_4 = 0 $）。因此，$ \{\boldsymbol{x}_1, \boldsymbol{x}_2, \boldsymbol{x}_4\} $ 是 $ U $ 的一组基。
\end{notebox}

\medskip\hrule\medskip

\section*{2.6.2 秩}
\addcontentsline{toc}{subsection}{2.6.2 秩}

矩阵 $ \boldsymbol{A} \in \mathbb{R}^{m \times n} $ 的线性无关列的数量等于其线性无关行的数量，并称为 $ \boldsymbol{A} $ 的 \emph{秩（rank）} ，记为 $ \operatorname{rk}(\boldsymbol{A}) $。

\emph{备注} 。矩阵的秩具有一些重要性质：

\begin{itemize}
\item $ \operatorname{rk}(\boldsymbol{A}) = \operatorname{rk}(\boldsymbol{A}^\top) $，即列秩等于行秩。
\item $ \boldsymbol{A} \in \mathbb{R}^{m \times n} $ 的列张成了 $ \mathbb{R}^m $ 的一个子空间 $ U \subseteq \mathbb{R}^m $，且 $ \operatorname{dim}(U) = \operatorname{rk}(\boldsymbol{A}) $。稍后我们将这个子空间称为 \emph{像（image）} 或 \emph{值域（range）} 。通过对 $ \boldsymbol{A} $ 执行高斯消元以确定主元列，可以求出 $ U $ 的一组基。
\item $ \boldsymbol{A} \in \mathbb{R}^{m \times n} $ 的行张成了 $ \mathbb{R}^n $ 的一个子空间 $ W \subseteq \mathbb{R}^n $，且 $ \operatorname{dim}(W) = \operatorname{rk}(\boldsymbol{A}) $。通过对 $ \boldsymbol{A}^\top $ 执行高斯消元，可以求出 $ W $ 的一组基。
\item 对于所有方阵 $ \boldsymbol{A} \in \mathbb{R}^{n \times n} $，$ \boldsymbol{A} $ 正则（可逆）当且仅当 $ \operatorname{rk}(\boldsymbol{A}) = n $。
\item 对于所有 $ \boldsymbol{A} \in \mathbb{R}^{m \times n} $ 和所有 $ \boldsymbol{b} \in \mathbb{R}^m $，线性方程组 $ \boldsymbol {Ax} = \boldsymbol{b} $ 有解当且仅当 $ \operatorname{rk}(\boldsymbol{A}) = \operatorname{rk}([\boldsymbol{A} \mid \boldsymbol{b}]) $，其中 $ [\boldsymbol{A} \mid \boldsymbol{b}] $ 表示增广矩阵。
\item 对于 $ \boldsymbol{A} \in \mathbb{R}^{m \times n} $，$ \boldsymbol {Ax} = \boldsymbol{0} $ 的解空间（子空间）的维度为 $ n - \operatorname{rk}(\boldsymbol{A}) $。稍后我们称这个子空间为 \emph{核（kernel）} 或 \emph{零空间（null space）} 。
\item 若矩阵 $ \boldsymbol{A} \in \mathbb{R}^{m \times n} $ 的秩等于相同维度矩阵所可能具有的最大秩，则称其为 \emph{满秩（full rank）} 。这意味着满秩矩阵的秩等于行数和列数中的较小者，即 $ \operatorname{rk}(\boldsymbol{A}) = \min(m, n) $。若矩阵不具有满秩，则称为 \emph{降秩（rank deficient）} 。
\end{itemize}

\begin{notebox}
\textbf{例 2.18（秩）}
\end{notebox}

\begin{notebox}
考察矩阵：

\begin{equation*}
\boldsymbol{A} = \begin{bmatrix} 1 & 0 & 1 \\ 0 & 1 & 1 \\ 0 & 0 & 0 \end{bmatrix}
\end{equation*}

$ \boldsymbol{A} $ 拥有两个线性无关的行/列，因此 $ \operatorname{rk}(\boldsymbol{A}) = 2 $。
\end{notebox}

\begin{notebox}
再考察矩阵：

\begin{equation*}
\boldsymbol{A} = \begin{bmatrix} 1 & 2 & 1 \\ -2 & -3 & 1 \\ 3 & 5 & 0 \end{bmatrix}
\end{equation*}

我们使用高斯消元法来确定其秩：

\begin{equation}
\begin{adjustbox}{max width=\linewidth}
\begin{bmatrix} 1 & 2 & 1 \\ -2 & -3 & 1 \\ 3 & 5 & 0 \end{bmatrix} \leadsto \cdots \leadsto \begin{bmatrix} 1 & 2 & 1 \\ 0 & 1 & 3 \\ 0 & 0 & 0 \end{bmatrix}
\end{adjustbox}
\tag{2.84}\label{eq:2-84}
\end{equation}

这里我们看到线性无关的行数和列数为 2，因此 $ \operatorname{rk}(\boldsymbol{A}) = 2 $。
\end{notebox}
